% Generated from validated Article Draft Result. Do not hand-edit; revise the structured draft instead.
\section{LLaMAからMixtralへ――open weightsは生態系になった}
\noindent\textbf{2023年は、公開されたweightsを起点に、規模、言語、専門化、効率、配布条件の異なるモデル系列が並走し始めた年だった。LLaMA/Llama 2、Falcon、Qwen、Mistral/Mixtralを通じて、open weightsをopen sourceやopen trainingと同一視せず、その多様化を追う。}\autocite{src-00c370ab2324,src-1fd43f93adcc,src-240ca564b6b8,src-3db09fede330,src-5c8e2db8871a,src-74aa7e2a7ca4,src-81f670ec1e74,src-8bdb904686fb,src-9df482e2c96c,src-b7e5243e6d59,src-bc7aba201f01,src-d4b6359bbad1,src-dd2cfa98160d,src-f0349196879e}

openという一語で括られがちなモデル群の公開範囲は同じではない。weightsが取得できること、利用条件が広いこと、学習データや学習手順まで再現可能であることは別の軸である。2023年のopen-weight拡大は、単一の公開モデルが現れたというより、異なる公開境界を持つ複数の系列が競争可能になった出来事として読む方が正確である。\autocite{src-00c370ab2324,src-1fd43f93adcc,src-240ca564b6b8,src-3db09fede330,src-5c8e2db8871a,src-74aa7e2a7ca4,src-81f670ec1e74,src-8bdb904686fb,src-9df482e2c96c,src-b7e5243e6d59,src-bc7aba201f01,src-d4b6359bbad1,src-dd2cfa98160d,src-f0349196879e}


モデルの技術的特徴だけでなく、licenseとaccessの条件が利用可能性を規定する。Llama 2以降の広い利用、FalconやQwenなど他系列の登場、Mistral系の公開は、weights公開の有無だけでは説明できない実装・再配布上の差を伴う。したがって本稿では「open」という形容ではなく、公開された成果物と利用条件を分けて記述する。\autocite{src-00c370ab2324,src-1fd43f93adcc,src-240ca564b6b8,src-3db09fede330,src-5c8e2db8871a,src-74aa7e2a7ca4,src-81f670ec1e74,src-8bdb904686fb,src-9df482e2c96c,src-b7e5243e6d59,src-bc7aba201f01,src-d4b6359bbad1,src-dd2cfa98160d,src-f0349196879e}


年の進行とともに、open-weightモデルは単発のbase modelからfamilyへ拡張した。汎用モデルに加え、instruction tuning、codeなどの専門化、複数サイズ、複数言語への展開が進み、利用者は一つの代表モデルを選ぶのではなく、用途と計算資源に合わせて系列内外を比較するようになった。\autocite{src-00c370ab2324,src-240ca564b6b8,src-3db09fede330,src-5c8e2db8871a,src-8bdb904686fb,src-9df482e2c96c,src-b7e5243e6d59,src-bc7aba201f01,src-d4b6359bbad1,src-dd2cfa98160d,src-f0349196879e}


年末のMixtralは、この多様化にarchitectureの選択肢をさらに加えた。ここで重要なのはMoEが既存方式を置き換えたと結論することではなく、公開weightsの世界でもparameter countだけでない効率設計が主要な比較軸へ入った点である。\autocite{src-00c370ab2324,src-dd2cfa98160d,src-f0349196879e}


\begin{claimboundary}
各プロジェクトや著者による性能比較は、その評価条件に基づく主張として扱う。公開範囲が広いことも、性能や再現性が自動的に保証されることを意味しない。\autocite{src-240ca564b6b8,src-3db09fede330,src-8bdb904686fb,src-9df482e2c96c,src-b7e5243e6d59,src-bc7aba201f01}
\end{claimboundary}
